\scanpage{247}
\mspage{193}

et en définissant symétriquement
\[
\Psi_\infty(X)=\underrightarrow{\lim}\bigl(X\to\Psi(X)\longrightarrow\unsure{\Phi|\mathscr{D}}\Psi(X)\longrightarrow\ldots\bigr).
\]
Pour que ce soit utilisable, il faut poser l'axiome

\uline{$W(6\ \mathrm{bis})$} (comparer $\overline{\underline{\mathrm{XII}}}$, pages $14,15$) La limite inductive d'une suite de morphismes $f_i:X_i\unsure{\mathbin{-}|\to}Y_i$ \unsure{i.e.|s.} pour un morphisme de systèmes inductifs dans $\mathrm{Cat}$ indexés par $\unsure{\mathbb N|N}$, est dans $W$ si les $f_i$ le sont.

Cet axiome est voisin de $W(6)$ de loc. cit., mais n'est impliqué par lui que moyennant $W(4\mathrm b)$). Il est plus fort dans le cas de syst. inductifs indexés par $N$ (ou aussi, plus gén., de syst. inductifs sur des ens. ordonnés filtrants ayant un ens. cofinal dénombrable), mais ne dit rien sur les syst. inductifs sur des ens. ordonnés d'indices plus généraux.

Cet axiome implique aussi qu'une limite inductive filtrante \deleted{dans $\underline{\mathrm{Cat}}_{/S}$} \inserted{de \unsure{morphismes|flèches}} \deleted{\unsure{est}} $f_i$ est $W$-lisse, resp. $W$-propre, si les $f_i$ le sont. \deleted{Ainsi, \unsure{pour}} On trouve alors que $\Phi_\infty(X)$ est à la fois $W$-lisse (comme limite inductive des $\Phi(X),\Phi\Psi\Phi(X)\ldots,$ $\Phi(\Psi\Phi)^n(X)$ \unsure{etc.}, \deleted{\unsure{cha}} \unsure{lesquelles|lesquels|quelles} sont $W$-lisses sur $S$) et $W$-propre (comme limite inductive des $\Psi\Phi(X),\ldots(\Psi\Phi)^n(X)\ldots$). \deleted{\unsure{Donc}} D'autre part, $X\longrightarrow\Phi_\infty(X)$ est dans $W_S^f$ -- car \inserted{\unsure{$W_S^f$ est stable par}} \deleted{\unsure{une}} limites inductives \inserted{filtr.} dénombrables dans $\mathrm{Cat}_{/S}$. On a donc trouvé un foncteur
% archive sheet number: 246
